\thispagestyle{plain}
\begin{titlepage}

\begin{center}

\LARGE{{\titel}}\\[1.5ex]
%\huge{\textbf{\titel}}\\[1.5ex]
%\LARGE{\textbf{\untertitel}}\\[6ex]
%\LARGE{\textbf{\art}}\\[1.5ex]
%\includegraphics[scale=0.7]{Hft_Logo.png}\\%\\[6ex]
%\includegraphics[scale=0.5]{Swiss_Post_Logo_inverted.png}%\\[6ex]

\end{center}
\footnotesize
\begin{tabular}{w{5cm}p{6cm}}\\\\\\
vorgelegt am:  &  12. Mai 2011\\[1.8ex]

Studienbereich: &  \studienbereich\\
Fakult\"at: &  \mbox{Vermessung, Informatik und Mathematik}\\
Bildungseinrichtung: &  \mbox{Hochschule f\"ur Technik Stuttgart}\\[1.8ex]


von:            &  \autor\\[1.8ex]

Matrikelnummer: &  \matrikelnr\\[1.8ex]

Pr\"ufer:  &  \mbox{\erstgutachter, HfT Stuttgart}\\[1.8ex]


\end{tabular}


\begin{center}
\textbf{Zusammenfassung}

\begin{abstract}
In Zeiten in denen davon gesprochen wird, dass Moore's Law in absehbarer Zukunft
nicht mehr gilt werden zusehens Architekturen eingesetzt, die nicht mehr darauf
vertrauen, dass jede einzelne CPU ausreichend schnell ist, sondern vielmehr 
massiv parallele Rechenleistung bieten. Im Zuge dieses Architekturwandels müssen
gleichwertig die eingesetzten Software-Stacks überdacht und teilweise neu
ausgelegt werden, um die neue Form von Rechenleistung konsolodiert und effizient
zu nutzen.

Clojure ist ein LISP Dialekt und damit eine primär funktionale
Programmiersprache, die es sich zum Ziel genommen hat dem Entwickler und der
Anwendung hohe Parallelität und Performanz bei gleichzeitig hohem Komfort
bereitzustellen. Diese Arbeit bietet einen Überblick über die dahinter stehende
Methodik und grenzt sie gegen den bisherigen Verfahren ab
\end{abstract}


%\copyright\ \jahr\\[9ex]

\end{center}

\end{titlepage}
